% Solution: Tadelis (2013) Exercise 4.7 "Campaigning". Input by \solitem.
% Reused from lista-reais-lectures-01-04-solutions (Tadelis 4.7), deepened. Check: check_A1.py.
\begin{exsol}{Tadelis Exercise 4.7}{77--78}{Exam}{3}
\asked Turn the campaigning story into a normal-form game and matrix, then run IESDS and say
whether it gives a clear prediction.

\method Normal-form game $G = (I, \{S_i\}, \{u_i\})$ (Lecture~1, slides~13--14; Tadelis \S3.1,
pp.~46--50); IESDS (Lecture~1, slides~28--29; Tadelis \S4.2, pp.~63--65). The only modelling step
is to fill the payoffs the story leaves implicit, using the fact that exactly one candidate wins.

\solution
\textbf{(a) The game.} Players $I = \{1, 2\}$; strategies $S_1 = S_2 = \{P, B, N\}$. Let
$\pi(s_i, s_j)$ be the probability that the candidate playing $s_i$ wins against one playing
$s_j$. The book gives four rules:
\[
  \pi(s,s) = 0.5,\qquad \pi(P,B) = 0,\qquad \pi(P,N) = 0.3,\qquad \pi(N,B) = 0.6 .
\]
Exactly one candidate wins, so $\pi(s_i,s_j) + \pi(s_j,s_i) = 1$. This fills the other three
entries:
\[
  \pi(B,P) = 1 - \pi(P,B) = 1,\qquad \pi(N,P) = 1 - 0.3 = 0.7,\qquad \pi(B,N) = 1 - 0.6 = 0.4 .
\]
Payoffs: $u_1(s_1,s_2) = \pi(s_1,s_2)$ and $u_2(s_1,s_2) = \pi(s_2,s_1) = 1 - u_1(s_1,s_2)$.
Player~2's payoff function is the same rule $\pi$ with the roles swapped (the game is
symmetric), and it is the complement of player~1's because win probabilities add up to one (the
game is constant-sum).

\textbf{(b) The matrix}, with best responses underlined:
\begin{center}
\begin{tabular}{r|c|c|c|}
\multicolumn{1}{r}{1 $\backslash$ 2} & \multicolumn{1}{c}{$P$} & \multicolumn{1}{c}{$B$} & \multicolumn{1}{c}{$N$}\\ \cline{2-4}
$P$ & $0.5,\ 0.5$ & $0,\ \br{1}$ & $0.3,\ \br{0.7}$ \\ \cline{2-4}
$B$ & $\br{1},\ 0$ & $0.5,\ 0.5$ & $0.4,\ \br{0.6}$ \\ \cline{2-4}
$N$ & $0.7,\ 0.3$ & $\br{0.6},\ 0.4$ & $\br{0.5},\ \br{0.5}$ \\ \cline{2-4}
\end{tabular}
\end{center}
$BR_1(P) = \{B\}$, $BR_1(B) = \{N\}$, $BR_1(N) = \{N\}$, and by symmetry
$BR_2(P) = \{B\}$, $BR_2(B) = \{N\}$, $BR_2(N) = \{N\}$.

\textbf{(c) IESDS.} Player~1's rows over columns $(P,B,N)$: $P = (0.5, 0, 0.3)$,
$B = (1, 0.5, 0.4)$, $N = (0.7, 0.6, 0.5)$.
\begin{itemize}
  \item \emph{Round 1.} $P$ vs $B$: $0.5 < 1$, $0 < 0.5$, $0.3 < 0.4$, so $P$ is strictly dominated
  by $B$. $P$ vs $N$: $0.5 < 0.7$, $0 < 0.6$, $0.3 < 0.5$, so $P$ is also dominated by $N$.
  $B$ vs $N$: $1 > 0.7$ but $0.5 < 0.6$, not ranked. By symmetry player~2's columns satisfy the
  same comparisons. $S_i^1 = \{B, N\}$ for $i = 1, 2$.
  \item \emph{Round 2.} On the columns $\{B, N\}$: $B = (0.5, 0.4)$, $N = (0.6, 0.5)$. $0.5 < 0.6$
  and $0.4 < 0.5$, so $B$ is strictly dominated by $N$. The same for player~2.
  $S_i^2 = \{N\}$.
  \item \emph{Round 3.} One strategy each: nothing to compare.
\end{itemize}
$S^\infty = \{(N, N)\}$: both run negative campaigns and each wins with probability $0.5$. The
procedure gives a clear prediction, a unique profile, after two rounds. In the marking $(N,N)$ is
the only cell with both payoffs underlined, so it is also the unique Nash equilibrium (as it must
be: Lecture~2, Proposition~2, slide~17).

\interp $B$ survives round~1 only because it crushes a positive campaign. Once each candidate
knows the other will not run a purely positive campaign, going fully negative wins. The outcome
has the same win probabilities as $(P,P)$, and in a constant-sum game every profile is Pareto
optimal \emph{for the candidates}. If voters dislike negative campaigning, the equilibrium is a
prisoner's-dilemma outcome for them: the candidates end up exactly where they would have been with
positive campaigns, but with a nastier campaign.

\verify \texttt{check\_A1.py} builds the matrix from the four rules plus complementarity, checks
symmetry, that $P$ is dominated by both $B$ and $N$, that IESDS removes $P$ (round~1) and $B$
(round~2) for both players, leaving $(N, N)$, and that $(N,N)$ is the only pure Nash equilibrium.
Tadelis gives no answer.

\mistakes
\begin{itemize}
  \item Filling the missing cells by symmetry of \emph{payoffs} ($\pi(B,P) = \pi(P,B) = 0$) instead
  of by complementarity. A candidate who beats $P$ for sure must win with probability~1.
  \item Deleting $B$ in round~1. Against $P$ it is the unique best reply ($1 > 0.7$); it becomes
  dominated only after $P$ is gone.
  \item In (a), writing only the matrix. The question asks for players, strategy sets and payoff
  functions, and for why player~2's function differs ($u_2 = 1 - u_1$).
\end{itemize}

\answer{(a) $I = \{1,2\}$, $S_i = \{P,B,N\}$, $u_1 = \pi(s_1,s_2)$ from the rules plus
$\pi(s_i,s_j) + \pi(s_j,s_i) = 1$, and $u_2 = 1 - u_1$ (the same rule from candidate~2's side).
(b) Matrix above. (c) Round~1 deletes $P$ (dominated by $B$ and by $N$); round~2 deletes $B$
(dominated by $N$); unique prediction $(N, N)$, each wins with probability $0.5$.}
\end{exsol}
